\documentclass[12pt,a4paper]{article}
\usepackage[utf8]{inputenc}
\usepackage[french]{babel}
\usepackage[T1]{fontenc}
\usepackage{geometry}
\geometry{margin=2.5cm}
\usepackage{titlesec}
\usepackage{enumitem}
\usepackage{xcolor}
\usepackage{fancyhdr}
\usepackage{tikz}

% Configuration des titres
\titleformat{\section}{\LARGE\textbf{\color{blue!80!black}}\centering}{}{0em}{}[\vspace{1ex}\titlerule]
\titleformat{\subsection}{\Large\textbf{\color{blue!60!black}}}{}{0em}{\medskip}
\titleformat{\subsubsection}{\large\textbf}{}{0em}{\smallskip}

% Mise en page
\pagestyle{fancy}
\fancyhf{}
\lhead{Obscurus Anti-Cheat Solutions}
\rhead{Fiche de Poste H/F - 2026}
\cfoot{\thepage}

\begin{document}

% --- FICHE : MAËL (PDG / CTO) ---
\section{DIRECTEUR GÉNÉRAL \& TECHNIQUE (H/F)}
\textbf{Titulaire : Maël}

\subsection{Identification du Poste}
\begin{itemize}[label=\textbullet]
    \item \textbf{Département} : Direction Générale / Direction Technique (R\&D)
    \item \textbf{Liaisons hiérarchiques} : Rapport direct au Conseil d'Administration
    \item \textbf{Liaisons fonctionnelles} : Direction Commerciale, Direction Juridique, Équipes Engineering
    \item \textbf{Finalité} : Garantir la souveraineté technologique de la solution anti-cheat et assurer la pérennité économique de l'entreprise
\end{itemize}

\subsection{Missions et Activités Principales}

\subsubsection{Pilotage de la Stratégie et du Business (Pôle CEO)}
\begin{itemize}[label=\textbullet]
    \item Définition des objectifs annuels et pluriannuels de croissance de l'entreprise.
    \item Élaboration et suivi du budget R\&D et d'exploitation, incluant la validation des investissements en infrastructure Cloud et serveurs.
    \item Reporting financier et stratégique régulier auprès des actionnaires et investisseurs.
    \item Garantie de la conformité éthique et respect du RGPD pour les méthodes de scan mémoire et de télémétrie.
\end{itemize}

\subsubsection{Direction de l'Ingénierie Kernel (Pôle CTO)}
\begin{itemize}[label=\textbullet]
    \item Conception de l'architecture du driver Windows (Ring 0) et des mécanismes d'auto-protection comme l'Anti-debugging et l'Anti-VM.
    \item Revue de code critique pour prévenir toute instabilité système (BSOD) chez les clients.
    \item Arbitrage technologique sur les langages (C++, Rust, Go) et les frameworks de traitement de données massives.
    \item Supervision du cycle de vie des détections, de l'identification d'un nouveau cheat au déploiement des heuristiques.
\end{itemize}

\subsubsection{Management des Ressources Humaines}
\begin{itemize}[label=\textbullet]
    \item Définition des besoins en recrutement pour les ingénieurs système et Data Scientists.
    \item Réalisation des entretiens annuels des directeurs techniques et leads de pôles.
    \item Maintien d'un haut niveau d'expertise interne via la formation continue sur les nouvelles vulnérabilités OS.
\end{itemize}

\subsection{Compétences et Aptitudes Requises}
\begin{itemize}[label=\textbullet]
    \item \textbf{Savoir-faire (Expertise)} : Architecture système (Windows Internals), gestion de la mémoire et cybersécurité offensive.
    \item \textbf{Savoir-être (Comportement)} : Leadership d'opinion reconnu par l'industrie, sang-froid en cas de crise technique majeure et vision à long terme.
\end{itemize}

\subsection{Critères d'Évaluation (KPIs)}
\begin{itemize}[label=\textbullet]
    \item Taux de détection (Detection Rate) des cheats connus et stabilité logicielle (absence de crashs/BSOD).
    \item Performance business mesurée par l'atteinte des objectifs de chiffre d'affaires et de marge brute.
\end{itemize}

\newpage

% --- FICHE : JOAN (COMMERCIAL & COM) ---
\section{CONSEILLER COMMERCIAL \& CHARGÉ DE COMMUNICATION (H/F)}
\textbf{Titulaire : Joan}

\subsection{Identification du Poste}
\begin{itemize}[label=\textbullet]
    \item \textbf{Département} : Direction Commerciale / Marketing
    \item \textbf{Finalité} : Convertir les studios de jeux vidéo et valoriser l'expertise technique de l'entreprise.
\end{itemize}

\subsection{Missions et Activités Principales}

\subsubsection{Développement Commercial et Prospection (B2B)}
\begin{itemize}[label=\textbullet]
    \item Chasse de comptes auprès des studios en phase de développement sur PC et Consoles.
    \item Qualification des besoins selon le type de jeu pour proposer la couche de protection adaptée (Kernel-level, Server-side, hybride).
    \item Présentation du Dashboard et négociation contractuelle basée sur le nombre d'utilisateurs actifs mensuels (MAU).
    \item Suivi d'intégration du SDK et vente de modules complémentaires (IA, anti-DDOS).
\end{itemize}

\subsubsection{Communication Institutionnelle et B2B}
\begin{itemize}[label=\textbullet]
    \item Création de supports d'aide à la vente, études de cas et livres blancs techniques.
    \item Organisation de la présence de l'entreprise sur les salons professionnels comme la GDC ou la Gamescom.
    \item Gestion du site web pour mettre en avant le positionnement expertise Kernel/Cloud.
\end{itemize}

\subsubsection{Communication et Gestion de Crise}
\begin{itemize}[label=\textbullet]
    \item Surveillance en temps réel des signaux faibles sur les réseaux sociaux (Reddit, Discord) pour identifier une faille de sécurité.
    \item Coordination de la "War Room" avec le CTO en cas de bypass massif pour aligner les éléments de langage.
    \item Rédaction de communiqués officiels gradués selon la gravité de l'incident technique.
    \item Interface avec les Community Managers des clients pour fournir des réponses unifiées aux joueurs.
\end{itemize}

\subsection{Compétences et Aptitudes Requises}
\begin{itemize}[label=\textbullet]
    \item \textbf{Savoir-faire (Expertise)} : Vente complexe (Solution Selling), culture gaming technologique (hooks, drivers) et maîtrise du CRM.
    \item \textbf{Communication} : Qualités rédactionnelles irréprochables et anglais courant de négociation.
    \item \textbf{Savoir-être (Comportement)} : Diplomatie face aux communautés hostiles, réactivité et persévérance commerciale.
\end{itemize}

\subsection{Critères d'Évaluation (KPIs)}
\begin{itemize}[label=\textbullet]
    \item Chiffre d'Affaires (New Business) signé et taux de conversion des prospects.
    \item E-réputation globale de la solution et sentiment de la communauté sur l'efficacité de l'outil.
\end{itemize}

\newpage

% --- FICHE : CAMILLE (RH & COMPTA) ---
\section{DIRECTEUR RH \& CHARGÉ DE COMPTABILITÉ (H/F)}
\textbf{Titulaire : Camille}

\subsection{Identification du Poste}
\begin{itemize}[label=\textbullet]
    \item \textbf{Département} : Ressources Humaines / Direction Financière
    \item \textbf{Finalité} : Recruter des experts rares et garantir la fiabilité des informations financières.
\end{itemize}

\subsection{Missions et Activités Principales}

\subsubsection{Stratégie de Recrutement d'Élite}
\begin{itemize}[label=\textbullet]
    \item Sourcing spécifique de profils atypiques (Kernel, Reverse Engineering) sur GitHub et forums de sécurité.
    \item Mise en place de tests techniques de haut niveau en collaboration avec le CTO.
    \item Valorisation de la marque employeur pour attirer les passionnés de bas-niveau et de lutte contre la fraude.
\end{itemize}

\subsubsection{Gestion de la Performance et QVT}
\begin{itemize}[label=\textbullet]
    \item Pilotage des packages salariaux (fixe, variable, actions) et suivi des plans de formation continue.
    \item Prévention du burn-out des équipes techniques lors des périodes de crises de bypass.
    \item Veille au respect du Code du Travail et rédaction des clauses de confidentialité (NDA) renforcées.
\end{itemize}

\subsubsection{Comptabilité Générale et Clients}
\begin{itemize}[label=\textbullet]
    \item Enregistrement des flux bancaires et gestion des factures fournisseurs (Cloud, serveurs, matériel de test).
    \item Émission des factures SaaS pour les studios selon le nombre de joueurs actifs (MAU).
    \item Suivi des règlements internationaux et gestion des relances clients en coordination avec les commerciaux.
    \item Établissement des déclarations fiscales (TVA, liasse fiscale) et support pour la paie.
\end{itemize}

\subsubsection{Reporting et Fiscalité R\&D}
\begin{itemize}[label=\textbullet]
    \item Centralisation des dépenses R\&D pour le montage des dossiers de Crédit Impôt Recherche (CIR/CII).
    \item Actualisation du plan de trésorerie prévisionnel et production des indicateurs de burn rate.
\end{itemize}

\subsection{Compétences et Aptitudes Requises}
\begin{itemize}[label=\textbullet]
    \item \textbf{Savoir-faire (Expertise)} : Chasse de tête, droit social, normes comptables (PCG) et Excel avancé.
    \item \textbf{Savoir-être (Comportement)} : Rigueur extrême, intégrité totale sur les données sensibles et intelligence émotionnelle.
\end{itemize}

\subsection{Critères d'Évaluation (KPIs)}
\begin{itemize}[label=\textbullet]
    \item Délai de recrutement (Time-to-hire) des experts et taux de rétention des talents clés.
    \item Fiabilité des comptes et respect des échéances fiscales et sociales.
\end{itemize}

\newpage

% --- FICHE : ANNA (RSO) ---
\section{RESPONSABLE DE LA SÉCURITÉ OFFENSIVE (H/F)}
\textbf{Titulaire : Anna}

\subsection{Identification du Poste}
\begin{itemize}[label=\textbullet]
    \item \textbf{Département} : Direction Technique (R\&D) / Pôle Sécurité
    \item \textbf{Liaisons hiérarchiques} : Rapport direct au CTO
    \item \textbf{Finalité} : Identifier les vulnérabilités de l'anti-cheat et valider la robustesse des défenses avant déploiement.
\end{itemize}

\subsection{Missions et Activités Principales}

\subsubsection{Red Teaming et Simulation d'Attaques}
\begin{itemize}[label=\textbullet]
    \item Développement de PoC (Proof of Concept) d'outils de triche internes (aimbots, wallhacks) pour tester les détections.
    \item Tentatives actives de contournement du driver Kernel (Ring 0) via des techniques de Manual Mapping ou Kernel ROP.
    \item Tests de résistance de l'infrastructure de télémétrie contre les injections de données faussées.
\end{itemize}

\subsubsection{Reverse Engineering et Veille Adverse}
\begin{itemize}[label=\textbullet]
    \item Analyse approfondie des cheats payants du marché pour comprendre leurs méthodes d'obfuscation.
    \item Audits de sécurité sur le code source de l'entreprise pour identifier des failles potentielles (buffer overflows).
\end{itemize}

\subsubsection{Collaboration et Remédiation (Purple Teaming)}
\begin{itemize}[label=\textbullet]
    \item Rédaction de rapports d'infiltration documentant chaque vulnérabilité avec score de criticité et preuve d'exploitation.
    \item Collaboration avec les développeurs Kernel pour concevoir des contre-mesures et renforcer l'auto-protection du logiciel.
\end{itemize}

\subsection{Compétences et Aptitudes Requises}
\begin{itemize}[label=\textbullet]
    \item \textbf{Savoir-faire (Expertise)} : Maîtrise experte d'IDA Pro, Ghidra et du debugging noyau (WinDbg).
    \item \textbf{Développement} : Capacité à coder des outils offensifs en C++, Assembly (x64) et Python.
    \item \textbf{Savoir-être (Comportement)} : Curiosité insatiable, éthique irréprochable et esprit critique sur la solidité des défenses.
\end{itemize}

\subsection{Critères d'Évaluation (KPIs)}
\begin{itemize}[label=\textbullet]
    \item Nombre de failles critiques découvertes en interne avant exploitation externe.
    \item Délai de reproduction d'un bypass signalé par la communauté ou détecté sur le marché.
    \item Qualité des préconisations de renforcement effectif du produit.
\end{itemize}

% --- BLOC À INSÉRER APRÈS LA PREMIÈRE FICHE (MAËL) ---
\newpage
\section*{ORGANIGRAMME DE LA STRUCTURE}
\vspace{3cm}
\begin{center}
\begin{tikzpicture}[
    node distance=1.5cm and 1cm,
    base/.style={rectangle, draw, rounded corners, fill=blue!5, text width=4.5cm, align=center, minimum height=1.2cm, font=\small},
    dir/.style={base, fill=blue!20, text width=5.5cm},
    tech/.style={base, fill=red!10}
]

    % Niveaux
    \node (mael) [dir] {\textbf{Maël} \\ Directeur Général \& CTO};
    
    \node (anna) [tech, below=2cm of mael] {\textbf{Anna} \\ Responsable Sécurité Offensive};
    
    \node (joan) [base, below left=3cm and 0.5cm of anna] {\textbf{Joan} \\ Conseiller Commercial};
    
    \node (camille) [base, below right=3cm and 0.5cm of anna] {\textbf{Camille} \\ Dir. RH \& Comptabilité};

    % Liens hiérarchiques (Maël vers les autres)
    \draw [-{Stealth}, thick] (mael) -- (anna);
    \draw [-{Stealth}, thick] (mael) -| (joan);
    \draw [-{Stealth}, thick] (mael) -| (camille);

    % Liens fonctionnels / collaboration avec étiquettes décalées vers le bas
    \draw [dashed, gray] (anna) -- (joan) 
        node[pos=0.7, xshift=-0.5cm, sloped, font=\tiny, black] {Alerte Crise};
        
    \draw [dashed, gray] (anna) -- (camille) 
        node[pos=0.7, xshift=0.5cm, sloped, font=\tiny, black] {Recrutement Expert};

\end{tikzpicture}
\end{center}
\vfill

\end{document}